% GENERATED FILE. Edit canonical lessons, then run npm run generate:books.
% canonical-source-hash: fnv1a64:4a0233098e691b1e
% canonical-lessons: TE-C120-samudram, TE-C120-sarassu, TE-C120-adavi, TE-C120-aku, TE-C120-matti

\chapter{Sea, Lake, Forest, Leaf, Soil}
\label{ch:te-sea-lake-forest-leaf-soil}
\hlchaptermodality{120}

\begin{chapteropening}
\textbf{By the end of this chapter:} \emph{I can say the sea, a lake, a forest, a leaf and soil.}

Complete the last lesson of chapter 120: i can say the sea, a lake, a forest, a leaf and soil.
\end{chapteropening}

\section[\texorpdfstring{\te{సముద్రం} (samudraṁ)}{సముద్రం (samudraṁ)}]{\te{సముద్రం} (samudraṁ) — the sea}

\label{lesson:TE-C120-samudram}



\begin{quote}
Before the new one: say the Telugu for a vada, then the Telugu for coffee.
\end{quote}

\subsection*{You'll want to know: \te{సముద్రం}}
\textbf{\te{సముద్రం}} — \emph{samudraṁ} — \textquotedblleft{}the sea\textquotedblright{}.

\begin{grammarlens}[title={What you've built}]
One more word for this chapter.
\end{grammarlens}

\subsection*{Your turn}
\begin{itemize}
\raggedright
  \item \emph{Say it:} \emph{samudraṁ}
  \item \emph{Say it:} \emph{samudraṁ}, once more
  \item \emph{Say it:} say \emph{samudraṁ}
  \item \emph{Recall:} say the Telugu for a vada, then the Telugu for coffee, then say \emph{samudraṁ} again
\end{itemize}

\begin{tcolorbox}[breakable,colback=teal!4,colframe=teal!35!black,title={Before you move on}]
What does \te{సముద్రం} mean? (\textquotedblleft{}The sea\textquotedblright{}.) Say it once more.
\end{tcolorbox}
\section[\texorpdfstring{\te{సరస్సు} (sarassu)}{సరస్సు (sarassu)}]{\te{సరస్సు} (sarassu) — a lake}

\label{lesson:TE-C120-sarassu}



\begin{quote}
Before the new one: say the Telugu for coffee, then the Telugu for the sea.
\end{quote}

\subsection*{You'll want to know: \te{సరస్సు}}
\textbf{\te{సరస్సు}} — \emph{sarassu} — \textquotedblleft{}a lake\textquotedblright{}.

\begin{grammarlens}[title={What you've built}]
One more word for this chapter.
\end{grammarlens}

\subsection*{Your turn}
\begin{itemize}
\raggedright
  \item \emph{Say it:} \emph{sarassu}
  \item \emph{Say it:} \emph{sarassu}, once more
  \item \emph{Say it:} say \emph{sarassu}
  \item \emph{Recall:} say the Telugu for coffee, then the Telugu for the sea, then say \emph{sarassu} again
\end{itemize}

\begin{tcolorbox}[breakable,colback=teal!4,colframe=teal!35!black,title={Before you move on}]
What does \te{సరస్సు} mean? (\textquotedblleft{}A lake\textquotedblright{}.) Say it once more.
\end{tcolorbox}
\section[\texorpdfstring{\te{అడవి} (aḍavi)}{అడవి (aḍavi)}]{\te{అడవి} (aḍavi) — a forest}

\label{lesson:TE-C120-adavi}



\begin{quote}
Before the new one: say the Telugu for the sea, then the Telugu for a lake.
\end{quote}

\subsection*{You'll want to know: \te{అడవి}}
\textbf{\te{అడవి}} — \emph{aḍavi} — \textquotedblleft{}a forest\textquotedblright{}.

\begin{grammarlens}[title={What you've built}]
One more word for this chapter.
\end{grammarlens}

\subsection*{Your turn}
\begin{itemize}
\raggedright
  \item \emph{Say it:} \emph{aḍavi}
  \item \emph{Say it:} \emph{aḍavi}, once more
  \item \emph{Say it:} say \emph{aḍavi}
  \item \emph{Recall:} say the Telugu for the sea, then the Telugu for a lake, then say \emph{aḍavi} again
\end{itemize}

\begin{tcolorbox}[breakable,colback=teal!4,colframe=teal!35!black,title={Before you move on}]
What does \te{అడవి} mean? (\textquotedblleft{}A forest\textquotedblright{}.) Say it once more.
\end{tcolorbox}
\section[\texorpdfstring{\te{ఆకు} (āku)}{ఆకు (āku)}]{\te{ఆకు} (āku) — a leaf}

\label{lesson:TE-C120-aku}



\begin{quote}
Before the new one: say the Telugu for a lake, then the Telugu for a forest.
\end{quote}

\subsection*{You'll want to know: \te{ఆకు}}
\textbf{\te{ఆకు}} — \emph{āku} — \textquotedblleft{}a leaf\textquotedblright{}.

\begin{grammarlens}[title={What you've built}]
One more word for this chapter.
\end{grammarlens}

\subsection*{Your turn}
\begin{itemize}
\raggedright
  \item \emph{Say it:} \emph{āku}
  \item \emph{Say it:} \emph{āku}, once more
  \item \emph{Say it:} say \emph{āku}
  \item \emph{Recall:} say the Telugu for a lake, then the Telugu for a forest, then say \emph{āku} again
\end{itemize}

\begin{tcolorbox}[breakable,colback=teal!4,colframe=teal!35!black,title={Before you move on}]
What does \te{ఆకు} mean? (\textquotedblleft{}A leaf\textquotedblright{}.) Say it once more.
\end{tcolorbox}
\section[\texorpdfstring{\te{మట్టి} (maṭṭi)}{మట్టి (maṭṭi)}]{\te{మట్టి} (maṭṭi) — soil, earth}

\label{lesson:TE-C120-matti}



\begin{quote}
Before the new one: say the Telugu for a forest, then the Telugu for a leaf.
\end{quote}

\subsection*{You'll want to know: \te{మట్టి}}
\textbf{\te{మట్టి}} — \emph{maṭṭi} — \textquotedblleft{}soil, earth\textquotedblright{}.

\begin{grammarlens}[title={What you've built}]
One more word for this chapter.
\end{grammarlens}

\subsection*{Your turn}
\begin{itemize}
\raggedright
  \item \emph{Say it:} \emph{maṭṭi}
  \item \emph{Say it:} \emph{maṭṭi}, once more
  \item \emph{Say it:} say \emph{maṭṭi}
  \item \emph{Recall:} say the Telugu for a forest, then the Telugu for a leaf, then say \emph{maṭṭi} again
\end{itemize}

\begin{tcolorbox}[breakable,colback=teal!4,colframe=teal!35!black,title={Before you move on}]
What does \te{మట్టి} mean? (\textquotedblleft{}Soil, earth\textquotedblright{}.) Say it once more.
\end{tcolorbox}
